\documentclass{article}
\usepackage[utf8]{inputenc}
\title{Milestone 4: RISC-V Vector Assembly Implementation}
\author{Kalman Filter Project}
\date{\today}

\begin{document}
\maketitle

\section{Introduction}
This report details the implementation, verification, and performance analysis of the Linear Kalman Filter (LKF) and Extended Kalman Filter (EKF) vectorised using the RISC-V Vector (RVV) extension.

\section{Function-by-Function Walkthrough}
The vectorised implementations leverage the RISC-V RVV 1.0 extension to accelerate the core matrix operations. The vector length is configured dynamically using \texttt{vsetvli}.

\subsection{Matrix Multiplication (\texttt{mat\_mul\_vec})}
The matrix multiplication kernel computes $C = A \times B$. The loop order is chosen to write the output matrix $C$ exactly once per output chunk to minimise memory traffic.
\begin{verbatim}
# For k = 0..K-1:
fld     ft0, 0(a6)                # scalar A[i][k]
vle64.v  v4, (a7)                 # B[k][j..j+VL-1]
vfmacc.vf v0, ft0, v4            # v0 += A[i][k] * B[k][j..j+VL-1]
\end{verbatim}

\subsection{Element-wise Operations}
Addition, subtraction, and scaled addition treat the matrices as flat arrays and use \texttt{vfadd.vv}, \texttt{vfsub.vv}, and \texttt{vfmacc.vf}.
\begin{verbatim}
vle64.v  v0, (a1)       # load A chunk
vle64.v  v4, (a2)       # load B chunk
vfadd.vv v0, v0, v4    # v0 = A + B
vse64.v  v0, (a0)       # store to C
\end{verbatim}

\subsection{Matrix-Vector Multiplication (\texttt{mat\_vec\_mul\_vec})}
This function uses an ordered reduction (\texttt{vfredosum.vs}) for deterministic and reproducible summation across runs.
\begin{verbatim}
vfmul.vv v8, v0, v4              # element-wise product
vfredosum.vs v12, v8, v12         # v12[0] = sum(v8)
\end{verbatim}

\subsection{Matrix Transposition (\texttt{mat\_transpose\_vec})}
Transposition is vectorised by column using strided loads (\texttt{vlse64.v}) and unit-stride stores (\texttt{vse64.v}).
\begin{verbatim}
vlse64.v v0, (a6), s4             # strided load, stride = N*8
vse64.v  v0, (a7)                  # unit-stride store
\end{verbatim}

\section{Numerical Verification}
The vectorised output is compared against the scalar assembly implementation to ensure numerical accuracy ($\epsilon \le 10^{-9}$).

\subsection{Average Absolute Error per Joint}
\begin{table}[ht]
\centering
\begin{tabular}{|l|r|r|}
\hline
Joint & LKF Vector Error & EKF Vector Error \\
\hline
pelvis & 2.45e-16 & 6.73e-18 \\
L5 & 1.19e-16 & 8.91e-17 \\
L3 & 5.14e-17 & 1.36e-16 \\
T12 & 9.27e-17 & 8.88e-17 \\
T8 & 1.79e-17 & 1.53e-16 \\
neck & 2.26e-16 & 1.92e-16 \\
head & 8.05e-17 & 5.75e-17 \\
shoulderRight & 9.65e-17 & 1.29e-16 \\
upperArmRight & 2.04e-16 & 8.24e-17 \\
forearmRight & 3.49e-16 & 8.12e-17 \\
handRight & 1.47e-16 & 4.50e-16 \\
shoulderLeft & 1.45e-16 & 2.61e-16 \\
upperArmLeft & 5.07e-17 & 3.36e-16 \\
forearmLeft & 9.28e-17 & 9.28e-17 \\
handLeft & 8.21e-17 & 2.15e-16 \\
upperLegRight & 1.36e-16 & 2.88e-16 \\
lowerLegRight & 7.12e-17 & 6.06e-18 \\
footRight & 2.51e-17 & 1.16e-17 \\
toeRight & 1.20e-16 & 1.86e-17 \\
upperLegLeft & 2.52e-16 & 1.55e-16 \\
lowerLegLeft & 1.59e-16 & 1.81e-16 \\
footLeft & 8.70e-17 & 3.60e-16 \\
toeLeft & 4.06e-17 & 9.09e-17 \\
\hline
\end{tabular}
\caption{Average absolute error per joint for vectorised LKF and EKF.}
\end{table}

\subsection{Average Absolute Error per State Component}
\begin{table}[ht]
\centering
\begin{tabular}{|l|r|r|}
\hline
Component & LKF Vector Error & EKF Vector Error \\
\hline
px & 1.54e-16 & 1.78e-16 \\
vx & 5.50e-17 & 3.15e-16 \\
ax & 3.41e-18 & 2.07e-17 \\
jx & 1.34e-19 & 8.56e-19 \\
py & 3.09e-16 & 4.33e-16 \\
vy & 4.25e-16 & 7.26e-16 \\
ay & 2.57e-17 & 4.52e-17 \\
jy & 1.00e-18 & 1.80e-18 \\
pz & 1.17e-17 & 4.07e-17 \\
vz & 4.89e-16 & 5.26e-17 \\
az & 3.25e-17 & 3.55e-18 \\
jz & 1.20e-18 & 1.36e-19 \\
\hline
\end{tabular}
\caption{Average absolute error per component for vectorised LKF and EKF.}
\end{table}

\subsection{Overall Error Bounds}
\begin{itemize}
\item LKF Max Error: 4.22e-15
\item LKF Min Error: 0.00e+00
\item EKF Max Error: 6.66e-15
\item EKF Min Error: 0.00e+00
\end{itemize}
All errors are well within the required $10^{-9}$ tolerance.

\subsection{Direct Comparison: Scalar vs Vector}
\begin{table}[ht]
\centering
\begin{tabular}{|l|r|r|}
\hline
Joint/Component & M3 Scalar Error & M4 Vector Error \\
\hline
LKF pelvis & 2.33e-16 & 2.45e-16 \\
LKF px & 1.54e-16 & 1.54e-16 \\
EKF head & 7.18e-17 & 5.75e-17 \\
EKF vy & 5.03e-16 & 7.26e-16 \\
\hline
\end{tabular}
\caption{Comparison of errors showing vectorisation did not degrade numerical accuracy.}
\end{table}

\section{Performance Analysis}
The speedup analysis compares the scalar (M3) code against the vectorised (M4) code.

\subsection{Instruction Count Reduction}
The instruction count was measured using the QEMU plugin \texttt{libinsn.so}.

\begin{table}[ht]
\centering
\begin{tabular}{|l|r|r|r|}
\hline
Filter & M3 Scalar Insns & M4 Vector Insns & Reduction Ratio \\
\hline
LKF & 6,665,955,583 & 3,547,738,380 & 1.88x \\
EKF & 6,668,931,253 & 3,550,632,658 & 1.88x \\
\hline
\end{tabular}
\caption{Instruction count reduction from M3 scalar to M4 vector.}
\end{table}

\subsection{Speedup Measurement}
The wall-clock time was measured over multiple frames using the \texttt{perf\_compare} utility running in QEMU. Note that under QEMU emulation, vector operations are simulated using scalar loops, resulting in a wall-clock slowdown, however the instruction count correctly reflects the reduction in architectural instructions executed.

\begin{table}[ht]
\centering
\begin{tabular}{|l|r|r|r|}
\hline
Filter & M3 Scalar Time (s) & M4 Vector Time (s) & Speedup (QEMU emulated) \\
\hline
LKF & 25.84 & 116.04 & 0.22x \\
EKF & 24.55 & 106.91 & 0.23x \\
\hline
\end{tabular}
\caption{Execution time speedup comparison.}
\end{table}

\end{document}
